\chapter{Of Health, the Infirmary, and the Mind}
\label{art:health}

\articlenote{In which is set down the care of citizen health, the charge of the Infirmary and the Physicians who tend it, mandatory care for health and the prevention of disease, and the disposition of citizens determined unable to sustain themselves through labor.}

\section{The Charge and the Physicians}

The Department of the Infirmary is charged with the health of every citizen: the prevention of disease, the care and recovery of the sick and injured, the certification of fitness for trade and for labor, and the keeping of the records of every citizen's health that the Pact requires. The Infirmary is located on floors within the Mids and is run by the Physicians of the Infirmary, led by a Head Physician appointed and confirmed under the shadowing process of Article 13.

The Infirmary's charge extends to the disorders and afflictions of a citizen's mind no less than of the citizen's body: a Physician who examines a citizen troubled in mind keeps the same record, in the same confidence, and by the same procedure as for any other complaint, and no citizen is treated for such a trouble save by a Physician under this Article.

\begin{clauses}
\item A Physician of the Infirmary is licensed by the Head Physician to undertake medical care, and may administer medicines, set broken bones, suture wounds, treat a citizen's disorders of mind, and determine the cause of death under Article 2; no other citizen may undertake such care save where no Physician is available and immediate aid is required to preserve life.
\item The Infirmary shall keep records of every citizen's health that the Infirmary observes or is told of: inoculations, treatments received, injuries, illnesses, and such conditions as affect the citizen's fitness for certain labors. These records are kept in the Judicial archive under Article 6, and remain confidential, open to the citizen named and to Judicial only, save where this Pact elsewhere requires disclosure.
\item A citizen shall have the right to refuse medical treatment, save where the Pact expressly compels it (treatment of a communicable disease that threatens the silo; the mandatory health screening set out in Section 2; the examination and restraint of a citizen found incompetent to stand for trial under Article 6).
\end{clauses}

\section{Mandatory Care and the Keeping of Health}

Every citizen is required to maintain the citizen's own health and to make regular use of the Infirmary's resources for the examination, inoculation, and prevention of disease, as set down by the Head Physician and the Mayor. A citizen who does not maintain such care may be reassigned by the relevant Department head to simpler labor less demanding of health or strength, or to retirement if no lighter labor is available.

\begin{clauses}
\item Every citizen of majority shall submit, at intervals set by the Head Physician, to a comprehensive health screening: examination of the citizen's body, testing of the citizen's health records, and attestation of fitness or unfitness for the citizen's current trade. A citizen found unfit by a screening may be reassigned or medically restricted under Article 13.
\item A citizen who refuses or repeatedly avoids a mandatory health screening may be compelled to attend by the Sheriff, provided Judicial is notified; a citizen compelled by the Sheriff is entitled to a hearing before Judicial as to whether the screening was in fact mandatory and lawfully required.
\item The use of recreational facilities (exercise yards, baths, steam rooms, organized sports) is made available by the Infirmary and the Department of Supply for every citizen's health, and citizens are encouraged, but not compelled, to make use of them. Where a citizen's health is observed to be seriously declining, the Head Physician may recommend reassignment or restricted labor as a condition of continued health.
\item Citizens of majority shall maintain awareness of emergency-health procedures: the location of the nearest aid station, the treatment of burns and wounds, the recognition of communicable illness and the obligation to report it. Teaching of such matters is the charge of Article 15 (schooling) and the Infirmary itself.
\end{clauses}

\section{Communicable Illness and the Isolation of the Sick}

The Infirmary has the authority to isolate a citizen believed to have a communicable illness, to compel treatment, and to restrict the citizen's movement and contact with others, all in service to the silo's collective health. No citizen's individual preference overrides this necessity.

\begin{clauses}
\item A Physician who believes a citizen to be suffering from a communicable disease shall isolate that citizen, notify Judicial and the Sheriff, and undertake such treatment as the Physician judges necessary. The citizen is entitled to food, water, and care during isolation, but may not leave the isolation area without the Head Physician's written permission.
\item A citizen who violates the isolation ordered by a Physician, or who knowingly exposes others to a communicable illness, is answerable under Article 17 for endangerment.
\item The silo's water, air, waste disposal, and food-handling are monitored by the Infirmary in consultation with the Department of Mechanical under Article 9 and the Department of Supply under Article 11, that the spread of disease through the silo's systems be prevented.
\end{clauses}

\section{Death and the Determination of Cause}

When a citizen dies, a Physician of the Infirmary shall examine the body and determine the cause of death, entering the finding in the Infirmary records and reporting to Judicial under Article 2. This is the cornerstone of the silo's records and the Judicial archive's understanding of how citizens pass.

\begin{clauses}
\item A Physician who examines a body shall, to the best of the Physician's knowledge, determine and report the cause of death, entering the finding upon the Infirmary's own Form of Death --- the citizen's name and naming-day, the floor and hour the body was found, the examining Physician's own name, and the finding itself, set down in that order and no other --- including whether the death is believed to be by illness, injury, age, or by the citizen's own will outside the clean request under Article 18.
\item Where a citizen is believed to have caused the citizen's own death by means other than a clean request under Article 18, the Physician shall report this finding to Judicial upon the same Form, and Judicial shall investigate and determine whether the finding is accurate before closing the rolls under Article 2. This investigation is not delayed; Judicial shall make its determination within a season.
\item Where a citizen dies by the cleaning procedure under Article 18, the Form of Death this section otherwise requires is not completed, and a Physician need not be present; Information Technology under Article 8 shall certify instead, by its own record and in its own keeping, that the sensors were properly maintained and the procedure was carried out as ordained, and this certification is sufficient for Judicial to close the rolls, provided no evidence suggests the cleaning was sought other than by the citizen's own final request.
\item A Physician who falsifies a cause of death, or who fails to report a suspected self-killing to Judicial, is answerable under Article 17 among the gravest offenses.
\end{clauses}

\section{Retirement and the Care of the Aged and Infirm}

A citizen who can no longer perform the labor assigned to the citizen—whether by age, illness, or injury—may be retired from labor under Article 13 and shall be cared for by the silo. The Infirmary and the Department of Supply together manage the provision of housing, food, medicine, and comfort for the retired.

\begin{clauses}
\item A citizen may petition for retirement after reaching the age at which the Head Physician and the Mayor jointly determine that further labor is no longer expected or possible. A citizen may also petition for early retirement if a chronic illness or permanent injury makes the citizen's current trade impossible; such a petition shall be referred to Judicial for a ruling.
\item A retired citizen continues to receive the ration of a citizen under Article 11, plus additional provisions for medicine or care if the Infirmary judges it necessary. A retired citizen's housing is assigned under Article 14 and may not be revoked save by the Head of Supply in consultation with the Head Physician.
\item A retired citizen who recovers strength or fitness sufficient to undertake lighter labor may petition to return to work; if approved by the Head Physician and the relevant Department head, the citizen may reenter a trade and will no longer receive retirement provisions beyond the ordinary ration.
\end{clauses}

\section{Physicians' Judgment and At-Risk Referral}

A Physician of the Infirmary holds responsibility for the wellbeing of the citizens under their care. Where a Physician believes a citizen to be in immediate and critical danger of self-harm, the Physician may refer that citizen to Judicial for protective evaluation and care.

\begin{clauses}
\item A Physician who believes a citizen to be in immediate, critical danger of self-harm or deterioration may refer that citizen to Judicial for protective evaluation, using the Physician's own judgment. No Physician is obligated to refer citizens believed to be at risk of lesser danger or slow decline.
\item A citizen who fears the citizen is in danger of self-harm may petition Judicial directly for protective evaluation, and may request the Infirmary's care and support without waiting for a Physician's referral. A citizen who petitions Judicial under this clause is entitled to a hearing as to the necessity and form of any protective measures.
\item A Physician who observes in a citizen a persistent fixation upon some matter this Pact does not teach or require, carried to the neglect of the citizen's own health, household, or labor, may refer the citizen to Judicial for evaluation under this section, using the Physician's own judgment; no Physician is obligated to make such a referral. A referral under this clause is not itself a charge against the citizen, and creates no mark upon the citizen's record under Article 21 save the fact of the referral and its resolution.
\end{clauses}

\section{Fitness for Critical Roles and the Syndrome}

Certain positions carry responsibilities that demand full physical and mental acuity. The Infirmary and Judicial work together to ensure that a citizen holding a critical role under Article 13 is fit to do so, and to identify cases where a citizen's affliction or infirmity requires immediate attention and possible removal from the role.

\begin{clauses}
\item A Physician of the Infirmary may, at any time, refer a citizen holding a critical role to Judicial if the Physician believes the citizen's health—mental or physical—is compromised in a way that threatens the silo's safety or the citizen's ability to perform the role. This referral is confidential until Judicial determines whether cause exists to remove the citizen.
\item No citizen afflicted with the Syndrome—a degenerative neurological condition characterized by tremors, progressive cognitive decline, and nervous-system deterioration—may hold a critical role, and any citizen in a critical role diagnosed with the Syndrome shall immediately report this to the Head Physician, resign from the role, and undertake the treatment the Infirmary recommends.
\item The Infirmary shall maintain confidential records of every citizen's Syndrome status, made known to Judicial only where a citizen holds or seeks a critical role. Diagnosis of the Syndrome does not mark a citizen as unfit for ordinary labor; it marks only that the citizen may not undertake work bearing direct responsibility for the silo's welfare or the lives of other citizens.
\end{clauses}
